\section{Résultats}
\label{sec:resultats}

\subsection{Audit de saturation}

La figure~\ref{fig:audit} résume l'audit. Sur \ChiffreAuditN{} tirages Sobol
de la boîte v5, $\log_{10}k_p$ reste dans $[\ChiffreKpPcinq;\ChiffreKpPqs]$
(5\textsuperscript{e} et 95\textsuperscript{e} centiles, médiane
\ChiffreKpMed) et touche sa borne basse dans \ChiffreKpSat{} des cas. La
saturation de $79~\%$ rapportée par la campagne antérieure exigerait
$k_p<10^{-4}$, soit quatre ordres de grandeur sous ce que produit la
projection. Elle signale donc un défaut \emph{en amont} de l'écrêtage, et non
une borne trop serrée.

Deux composantes saturent à $100~\%$ : $\log_{10}\omega_f$ et $b$. Toutes
deux sont \emph{constantes} de la synthèse ($\omega_f=1/(r_CC)$ et $b=1$ par
la relation~\eqref{eq:projection}) : la saturation est structurelle et ne
retire aucun degré de liberté. Le vecteur FO-PIDF issu de la synthèse n'a
donc que cinq degrés de liberté effectifs, et non sept.

\begin{figure}[t]
\centering
\includegraphics[width=0.9\linewidth]{fig_audit_saturation.pdf}
\caption{Audit de saturation des bornes d'ancrage. Aucune composante variable
ne sature ; les deux saturations observées portent sur des constantes de la
synthèse.}
\label{fig:audit}
\end{figure}

\subsection{Sonde exploratoire : la région faisable existe}

Avant la campagne, une sonde exploratoire (graines 1 à 6, distinctes des
graines de campagne) a cherché conjointement sur tout $\theta$, à 1500 appels
par méthode. La région nominalement faisable \emph{existe} : avec $b$ libre,
la meilleure marge vaut \ChiffreMargeBLibre{} après rejeu à $0{,}2~\mu$s
(\ChiffreNFaisBLibre{} méthodes sur \ChiffreNLancersLibre{} l'atteignent) ;
avec $b=1$ imposé, elle vaut encore \ChiffreMargeBUn{}
(\ChiffreNFaisBUn{} sur \ChiffreNLancersUn), figure~\ref{fig:levierb}. Le
paramètre $b$ est donc un levier de marge et non une barrière de faisabilité,
contrairement à l'hypothèse que nous avions d'abord formulée. En revanche, le
point faisable trouvé avec $b=1$ exige $\log_{10}k_p\approx-1{,}66$, hors de
l'image de la boîte v5 par la synthèse ; celui trouvé avec $b$ libre exige
$b=0$, que v5 ne peut pas régler.

\begin{figure}[t]
\centering
\includegraphics[width=0.9\linewidth]{fig_levier_b.pdf}
\caption{Sonde exploratoire : meilleure marge nominale par méthode, $b$ libre
contre $b=1$. Recherche conjointe sur tout $\theta$, 1500 appels par lancer.}
\label{fig:levierb}
\end{figure}

\subsection{Campagne confirmatoire}

Le tableau~\ref{tab:bras} résume les trois bras. \textbf{Aucune des
540 cellules n'est robustement faisable sur le jeu de conception}, et
l'intervalle de Wilson par méthode et par bras est $[0;\ChiffreWilsonHaut]$.

\begin{table}[t]
\centering
\caption{Campagne confirmatoire, 30 graines $\times$ 6 méthodes par bras.}
\label{tab:bras}
\begin{tabular}{lccc}
\toprule
 & v5 (cadrage) & v5b ($b$ libre) & v5kp ($k_p$ libre) \\
\midrule
cellules & \ChiffreNCinq & \ChiffreNCinqB & \ChiffreNCinqKp \\
graines stériles (étage 1) & \ChiffreSterilesCinq & \ChiffreSterilesCinqB & \ChiffreSterilesCinqKp \\
cellules avec un point nominal faisable & \ChiffreNomFaisCinq & \ChiffreNomFaisCinqB & \ChiffreNomFaisCinqKp \\
meilleure marge nominale (étage 1) & \ChiffreMargeEtUnCinq & \ChiffreMargeEtUnCinqB & \ChiffreMargeEtUnCinqKp \\
cellules robustement faisables & \ChiffreRobFaisCinq & \ChiffreRobFaisCinqB & \ChiffreRobFaisCinqKp \\
meilleure marge robuste (étage 2) & \ChiffreMargeEtDeuxCinq & \ChiffreMargeEtDeuxCinqB & \ChiffreMargeEtDeuxCinqKp \\
Friedman, $p$ & \ChiffreFriedmanPCinq & \ChiffreFriedmanPCinqB & \ChiffreFriedmanPCinqKp \\
paires significatives (Holm, sur 15) & \ChiffrePairesCinq & \ChiffrePairesCinqB & \ChiffrePairesCinqKp \\
méthode de meilleure marge médiane & \ChiffreMeilleureCinq & \ChiffreMeilleureCinqB & \ChiffreMeilleureCinqKp \\
\bottomrule
\end{tabular}
\end{table}

Le bras v5, conforme au cadrage, ne trouve \emph{aucun} point nominalement
faisable en 700 appels, pour aucune des 180 cellules. Libérer $b$ (v5b) rend
la faisabilité nominale accessible dans \ChiffreNomFaisCinqB{} cellules sur
180, et libérer $k_p$ (v5kp) dans \ChiffreNomFaisCinqKp{} ; dans les deux cas,
aucune cellule ne tient sur le jeu de conception.

Les tests de Friedman sont significatifs dans chaque bras. Ces différences
portent toutefois sur la \emph{proximité} de la faisabilité (des marges
toutes négatives), et non sur une faisabilité atteinte. Le classement des
méthodes change d'un bras à l'autre (\ChiffreMeilleureCinq{} en tête sur v5,
\ChiffreMeilleureCinqB{} sur v5b) : aucune supériorité universelle n'est
établie (figure~\ref{fig:graines}).

\begin{figure}[t]
\centering
\includegraphics[width=0.49\linewidth]{fig_graines_v5.pdf}\hfill
\includegraphics[width=0.49\linewidth]{fig_graines_v5b.pdf}
\caption{Marge robuste par graine (unité statistique), bras v5 et v5b.
Échelle symlog : linéaire sur $[-1;1]$.}
\label{fig:graines}
\end{figure}

\paragraph{Relance après correction du feedforward.}
La campagne a d'abord été calculée avec un défaut : le feedforward de
l'étage~2 utilisait les paramètres du plant perturbé au lieu des valeurs
nominales figées. Elle a été relancée en entier après correction (protocole
v6.2), avec les mêmes graines, budgets et méthodes. Les chiffres rapportés
ici sont ceux de la relance. L'étage~1 se reproduit à l'identique dans les
540 cellules. À l'étage~2, les meilleures marges robustes passent de
\ChiffreAncMargeEtDeuxCinq{}, \ChiffreAncMargeEtDeuxCinqB{} et
\ChiffreAncMargeEtDeuxCinqKp{} à \ChiffreMargeEtDeuxCinq{},
\ChiffreMargeEtDeuxCinqB{} et \ChiffreMargeEtDeuxCinqKp{} (bras v5, v5b et
v5kp) ; aucune cellule ne devient faisable et le classement des méthodes est
inchangé.

\paragraph{Étage 3 et test scellé.}
Aucun candidat n'étant faisable en conception, l'étage~3 n'a rien à rejouer
et le test scellé \emph{n'a pas été joué}. Il reste vierge, graine
\ChiffreGraineTest{} comprise, et conserve toute sa valeur probante pour une
étude ultérieure.

\subsection{Validation croisée avec le modèle Simulink}

Six candidats figés ont été rejoués sur le modèle Simulink de référence (plant
nominal), avec le même correcteur, traduit en S-function, et les mêmes
formules de coût et de contraintes. Les critères d'accord avaient été fixés
avant l'essai. Les \ChiffreXvalN{} candidats dotés d'un correcteur sont en
accord : écart sur $J$ d'au plus \ChiffreXvalDJmax, écart sur chaque $g_i$
d'au plus \ChiffreXvalDgmax, verdict de faisabilité identique. Les deux points
faisables de la sonde restent faisables sous Simulink (marge
\ChiffreXvalMargeSl{} pour le point à $b=1$).

Le témoin sans rétroaction ne satisfait pas le critère sur $g_4$ (écart
\ChiffreXvalFFdg) : son erreur au palier de 1~V diffère de
\ChiffreXvalFFde{} entre les deux modèles. Sans rétroaction, rien ne corrige
l'équilibre en conduction discontinue, qui dépend du modèle de diode et
d'interrupteur ; en boucle fermée, l'intégrateur efface cet écart. La même
origine explique un $V_\mathrm{min}$ plus haut de \ChiffreXvalDVmin{} en
moyenne sous Simulink. Le plancher à $d_\mathrm{min}$ y reste sous le seuil
($V_\mathrm{min}=\ChiffreXvalVminSl$) : le verrou physique est confirmé par le
modèle de référence.

\subsection{Deux diagnostics post hoc}

Les deux diagnostics suivants ont été ajoutés \emph{après} observation. Ils
sont rapportés, mais n'ont servi à aucune sélection.

\paragraph{Plancher de tension à rapport cyclique minimal.}
Les marges $g_6$ reviennent à l'identique, à la quatrième décimale près,
d'un correcteur et d'un bras à l'autre. À $d=d_\mathrm{min}=0{,}02$ en boucle
ouverte, le convertisseur en conduction discontinue se stabilise à
\ChiffrePlancherNominal{} au nominal, et sous $0{,}8$~V dans
\ChiffrePlancherN{} des \ChiffrePlancherTot{} cas (jusqu'à
\ChiffrePlancherMin). Tout correcteur qui reste saturé au plancher pendant la
descente $24\to1$~V y est piégé et viole $g_6$ ; pour satisfaire $g_6$, il
doit se désaturer avant que la tension ne passe sous $0{,}8$~V.

\paragraph{Cycles limites invisibles aux contraintes.}
Le meilleur correcteur de la sonde, faisable au sens des six contraintes,
entretient au palier de 24~V un cycle limite à environ $7{,}7$~kHz, de
624~mV crête à crête sur $v_o$ moyennée par période MLI
(figure~\ref{fig:traj}). Le correcteur obtenu avec $b=1$ oscille aussi, plus
faiblement (158~mV). Les erreurs de palier sont des moyennes de fenêtre, qui
annulent l'oscillation, et le terme en $|\Delta d|$ de $J$ est trop faible pour
la pénaliser. Avec $b=0$, la sortie chute en outre d'environ 24 à 14~V au
démarrage, transitoire qu'aucune contrainte ne voit.

\begin{figure}[t]
\centering
\includegraphics[width=0.85\linewidth]{fig_trajectoire_b_libre.pdf}
\caption{Trajectoire du meilleur correcteur de la sonde ($b$ libre).
Faisable au sens des six contraintes, il présente un creux au démarrage et
un cycle limite au palier de 24~V.}
\label{fig:traj}
\end{figure}
